% Full local proofs for the finite-prefix normalization and Gold I.9.

\begin{annexitem}{annex:p15-pr-learning-coding}{thm:p15-pr-learning-coding}{Finite configurations, histories and bounded simulation}
\begin{proof}
One concrete finite-sequence coding uses the pairing function
\[
 \pi(u,v)=\frac{(u+v)(u+v+1)}2+v
\]
and codes the empty sequence by \(0\) and the sequence with first
entry \(u\) and remaining-sequence code \(v\) by
\(\pi(u,v)+1\). The two inverse pairing projections are
primitive recursive: search the finite square
\(0\leq u,v\leq z\) for the unique pair with \(\pi(u,v)=z\).
Every bounded search is implemented by primitive recursion,
with a default on an empty search.
The tail of a nonzero sequence code is strictly smaller than
that code. Iterating the tail at most the original code many
times therefore finds its length and entries by bounded
recursion. Sequence validity, truncation, appending and
prefix construction are primitive recursive too: extract
the boundedly many relevant entries and rebuild by the
same constructor. These operations need not have small
complexity bounds; their membership in the primitive-recursive
class is the point.

For deterministic finite-program machines, code a configuration
by its state and two finite stacks of tape symbols, with the
currently scanned symbol at the top of the right stack.
Use a fixed finite tape alphabet with blank \(0\) and a base \(B\)
larger than every symbol code. A stack is coded by its base-\(B\)
digits, with its top digit least significant. Reading and
popping use remainder and quotient by \(B\); pushing a
symbol \(d\) onto a stack \(s\) uses \(d+Bs\).
Popping the empty stack returns blank \(0\) and empty stack \(0\).
Leading blanks can be omitted because they represent the same
infinite blank tail. More explicitly, write the stacks as
\(l,r\), and put \(r_{\mathrm{tail}}=\lfloor r/B\rfloor\).
If the transition table on the scanned symbol \(r\bmod B\)
prescribes written symbol \(w\) and a move right, the new
stacks are \(w+Bl,r_{\mathrm{tail}}\).
For a move left they are
\[
 \left\lfloor l/B\right\rfloor,\qquad
 (l\bmod B)+B\bigl(w+B r_{\mathrm{tail}}\bigr).
\]
The new state is the table entry's prescribed state.
These formulas also cover an empty blank stack and keep
the head's newly scanned symbol at the right stack's top.
Finite state/symbol transition-table lookup is a bounded
search through the finite program description.
Quotients, remainders, finite tuple projections and those
bounded searches are primitive recursive. Thus initialization,
one transition and the halting-state predicate are primitive
recursive, uniformly on finite program descriptions.
Malformed descriptions or configurations receive a fixed
default transition; that total extension does not certify
them as valid histories.

Fix binary numerical input and output conventions, with
tape digits \(1,2\) for binary symbols \(0,1\), and digit
\(3\) for a delimiter; a base at least \(4\) suffices.
At input or numerical halting output, the head scans the
first digit of a nonempty finite digit block, followed
to the right by the delimiter. Use least-significant
binary digit first, and represent numerical \(0\) by
one binary-zero symbol. The output extractor verifies
this finite format and evaluates its binary digits;
trailing tape content beyond the delimiter is irrelevant.
Producing the finite input for \(x\) uses repeated
division by \(2\) for at most \(x+1\) rounds. It is primitive recursive.
Reading a halting configuration's finite output uses bounded
digit extraction and finite polynomial evaluation.
Arbitrary finite alphabet machines can use this fixed
alphabet by finite block coding; the interpretation \(Z_a\)
and its input/output convention are fixed once, rather
than changed during the diagonal.

Let \(y\) code a nonempty sequence of configurations.
Define \(T(a,x,y)\) to say that the description of \(a\)
is well formed, the first configuration is its initialization
on \(x\), each adjacent pair is a legal nonhalting transition,
and the last configuration is halted with a well-formed
numerical output. All checks range over the finite sequence
decoded from \(y\); bounded conjunction of primitive-recursive
conditions is primitive recursive. Define \(U(y)\) to extract
that final output, and give it value \(0\) on malformed codes.
Every actual finite halting computation has such a code.
Conversely a verified history follows the deterministic
machine from exactly the specified initialization, so its
last output equals \(Z_a(x)\). This proves (PR.1).
If two verified histories have different lengths because
of an alternative harmless stopping-code convention, they
still have the same numerical output.

Finally, if \(c_0=\operatorname{Init}(a,x)\) and
\(c_{s+1}=\operatorname{Step}(a,c_s)\), then
\((a,x,s)\mapsto c_s\) is primitive recursive by the
primitive-recursion clause itself. Halting configurations
may be made stationary for this bounded simulator.
Unbounded search for a halting stage is not part of this
construction and is not asserted to be total.
\end{proof}
\end{annexitem}

\begin{annexitem}{annex:p15-pr-learning-normalization}{thm:p15-pr-learning-normalization}{Serial completion preserves both convergence directions}
\begin{proof}
For fixed input \(\sigma\), code the simulator's state as
the tuple containing next-prefix number \(m\), active/inactive
flag, simulated machine configuration and retained answer.
The update at one tick uses only the operations in
\cref{thm:p15-pr-learning-coding}. Initialize an inactive call
on the extracted prefix, apply one transition, and check
whether the result is halted. If it is halted, record its
answer, increase \(m\) and reset the flag; otherwise retain
the active configuration. An initially halted computation
uses the stationary transition and still emits only one
answer on that tick. The initial state has \(m=1\).
Induction on \(r\) shows \(m\leq r\) immediately before
tick \(r\): at most one call completed at each earlier
tick. Thus every extraction at ticks \(r\leq n\)
asks for an available prefix.
This finite-state update is primitive recursive in
the input code and the old state. Iteration for
\(n=\operatorname{length}(\sigma)\) ticks and final
answer projection are primitive recursive.
On invalid sequence codes use a fixed default answer.
This proves total primitive recursiveness on all numerical
inputs as well as the asserted valid-prefix behavior.

Fix an infinite sequence \(i\). Running the procedure
on its prefix of length \(n+1\) agrees with the length-\(n\)
procedure on its first \(n\) ticks: at tick \(r\leq n\)
only a prefix of length \(m\leq r\) is requested, and
those entries are identical. The initial states and
deterministic transitions agree. Therefore the answers
\(\widehat G(i_1,\ldots,i_n)\) are successive retained
outputs of one compatible infinite serial simulation,
not unrelated computations on unrelated initial data.

Write \(a_m=G_0(i_1,\ldots,i_m)\).
Because \(G_0\) is total, its first prefix computation
halts after finitely many ticks. Once the \((m-1)\)-st
call has completed, the \(m\)-th is initialized on
its correct prefix and also halts after finitely many
further ticks. Induction gives a completion tick
\(s_m<\infty\) for every \(m\).
Since at most one call completes per tick,
\(s_1<s_2<\cdots\); in particular these ticks tend
to infinity. The retained output equals \(a_m\)
at tick \(s_m\) and until the next completion.
No \(a_m\) is omitted or replaced by an answer from
a later prefix.

If \(a_m=a\) for all \(m\geq M\), then after tick
\(s_M\) the retained output is always \(a\).
This is the forward implication in (PR.2).
Conversely suppose retained outputs equal \(a\)
at every tick \(n\geq N\).
Every sufficiently large completion tick \(s_m\)
lies beyond \(N\), and the output there is exactly
\(a_m\). Hence \(a_m=a\) for all sufficiently large
\(m\). This is the reverse implication.
It proves preservation of nonstabilization too.
No alteration of \(i\), no semantic equivalence check
between hypotheses, and no known bound on any
completion time have been used.
\end{proof}
\end{annexitem}

\begin{annexitem}{annex:p15-pr-learning-range}{thm:p15-pr-learning-range}{Witness enumeration with a fixed positive default}
\begin{proof}
Choose a fixed word number \(a_L\in L\), which exists
by nonemptiness. Let \(e\) be a recognizer that halts
with numerical output \(1\) exactly on the word numbers
in \(L\). Such a recognizer exists by computable
enumerability; on other inputs it may diverge.
The fixed history predicate
\[
 R(x,y)\quad\Longleftrightarrow\quad
 x>0\ \land\ T(e,x,y)\ \land\ U(y)=1
\]
is primitive recursive, and
\(x\in L\) iff some \(y\) satisfies \(R(x,y)\).
Decode \(n-1\) by the primitive-recursive inverse
of \(\pi\), writing the result as \((x_n,y_n)\).
Define
\[
 f(n)=
 \begin{cases}
 x_n,&R(x_n,y_n),\\
 a_L,&\text{otherwise},
 \end{cases}
 \qquad n\geq1.
 \tag{PR.4}
\]
Both branch values and the branch test are primitive
recursive. Every output belongs to \(L\), including
the default. For each \(x\in L\), some finite
accepting-history code \(y\) exists, and the pair
\((x,y)\) occurs in the decoding of
\(n=\pi(x,y)+1\). At that position \(f(n)=x\).
Therefore its range is exactly \(L\).
If primitive recursion is defined on all of \(\mathbb N\),
set \(f(0)=a_L\); the displayed positive-index map is the
restriction of that total primitive-recursive extension.

The fixed default may depend on \(L\); there is no
claim of an algorithm finding it for every possible
recognizer index, including empty targets.
This is the nonempty range construction used in
the source's choice of \(f\), not a uniform list
of total decision algorithms for all recursive languages.
\end{proof}
\end{annexitem}

\begin{annexitem}{annex:p15-pr-learning-obstruction}{thm:p15-pr-learning-obstruction}{Gold I.9 with its normalization interface and all cases}
\begin{proof}
Suppose \(G_0\) identifies the declared class
\(\mathcal C\) from every primitive-recursive pure text
with correct tester names. If its infinite member \(L\)
were nonrecursive, no correct total tester could be
the final hypothesis. Thus we may assume \(L\)
recursive, fix \(f\) as in
\cref{thm:p15-pr-learning-range}, and normalize \(G_0\)
by \cref{thm:p15-pr-learning-normalization}.
The resulting primitive-recursive rule \(G\)
has exactly the same stabilized hypothesis on
every presentation on which either rule stabilizes.
If it fails on a permitted text, \(G_0\) fails there too.

Construct \(X_t,i_t,g_t,P_t,Q_t\) by (PR.3).
The prefix and \(X_t\) are already available when
\(g_t\) is computed; \(P_t\) is a bounded disjunction
over that prefix, and \(Q_t\) a bounded existential
search through history numbers \(0,\ldots,t\).
Both are primitive recursive.
For an explicit recursion, maintain the state
\((X_t,s_t)\), where \(s_t\) codes
\((i_1,\ldots,i_t)\).
Compute \(g_t=G(s_t)\), the two predicates,
\(X_{t+1}\) and \(i_{t+1}\); then append
\(i_{t+1}\) to obtain \(s_{t+1}\).
This is ordinary primitive recursion on a coded
pair of numbers. Extracting the last entry shows
that \(t\mapsto i_t\) is primitive recursive.
Equivalently recurse from the initial state at offset \(t-1\);
an unused value \(i_0=f(1)\) gives a total extension to
\(\mathbb N\) when that domain convention is required.
There is no unbounded search for a rejecting computation.
By induction, each \(i_t\in L\):
the initial value is \(f(1)\), and a subsequent
value either repeats its predecessor or is \(f(X_t)\).
In particular the content is nonempty.

\emph{Case I: \(\neg P_t\) holds only finitely often.}
Choose \(\alpha\) such that \(P_t\) holds for all
\(t\geq\alpha\).
The recursion then gives
\[
 X_t=X_\alpha+t-\alpha,\qquad
 i_t=i_\alpha\quad(t\geq\alpha).
\]
The entire text range is the finite set
\(F=\{i_1,\ldots,i_\alpha\}\).
Since \(P_t\) holds, every \(f(X_t)\) for
\(t\geq\alpha\) belongs to \(F\).
These \(X_t\) exhaust the integers at least \(X_\alpha\).
Only finitely many earlier values of \(f\) remain;
its range would be finite, contrary to its range \(L\).
Thus this case cannot occur, and \(\neg P_t\)
holds infinitely often.

\emph{Case II: \(Q_t\) holds only finitely often.}
Choose \(\alpha\) such that \(Q_t\) is false for
all \(t\geq\alpha\). The update of \(i\) then
repeats \(i_\alpha\) forever, so the text range
is the finite nonempty language
\(F=\{i_1,\ldots,i_\alpha\}\in\mathcal C\).
If \(g_t\) does not stabilize, identification fails.
Otherwise let its eventual value be \(a\).
Because \(\neg P_t\) occurs infinitely often,
choose \(\beta\geq\alpha\) beyond stabilization
with \(\neg P_\beta\). Set \(x=f(X_\beta)\).
It is absent from \(F\).
From stage \(\beta\) onward the content remains
\(F\), \(X_t=X_\beta\) and \(P_t\) stays false:
this follows by induction from the fixed content
and the update rule for \(X\).
Since \(Q_t\) is false for every such \(t\),
\[
 \neg(\exists y\leq t)
 [T(a,x,y)\land U(y)=0]\quad(t\geq\beta).
\]
If \(Z_a(x)\) halted with output \(0\), (PR.1)
would supply one finite history number \(y_0\);
taking \(t\geq\max(\beta,y_0)\) contradicts this
display. Thus \(a\) does not reject \(x\):
it may diverge, return \(1\), or return another value.
Every correct tester for \(F\) must reject
the word \(x\notin F\) with output \(0\).
The stabilized hypothesis is therefore not correct.

\emph{Case III: \(Q_t\) holds infinitely often.}
We first prove that the text range is exactly \(L\).
The inclusion in \(L\) was already proved.
The sequence \(X_t\) starts at \(1\), is nondecreasing,
and increases only by \(1\). Suppose it is eventually
constant at \(X\). Then \(P_t\) must be false
at every sufficiently late stage, since a true
\(P_t\) would increase it. Choose a later stage
\(\beta\) with \(Q_\beta\). Then the insertion
rule puts \(i_{\beta+1}=f(X)\).
Consequently \(P_{\beta+1}\) is true and
\(X_{\beta+2}=X+1\), a contradiction.
Thus \(X_t\) is unbounded and visits every positive
integer \(k\). At its last stage \(t\) with
\(X_t=k\), it advances to \(k+1\); the advancement
rule forces \(P_t\), so \(f(k)\) has already
appeared in the text. Every value of \(f\)
therefore appears. This also makes explicit the
source's bounded/constant and passing subcases:
a missing \(f(k)\) can neither be held at forever
nor be passed without first having been observed.
The text is fair for \(L\), even though repetitions
and delays may be arbitrarily long.

If \(g_t\) fails to stabilize, the conclusion
again follows. If \(g_t=a\) for all
\(t\geq\alpha\), choose \(\beta\geq\alpha\)
with \(Q_\beta\). By its definition there is
a witnessed computation
\[
 T[a,f(X_\beta),y]\land U(y)=0
 \quad\text{for some }y\leq\beta.
\]
Hence \(Z_a(f(X_\beta))=0\) by (PR.1).
But \(f(X_\beta)\in L\), the text's full range.
A correct tester for \(L\) must answer \(1\) there.
The stabilized hypothesis is wrong.
It is immaterial whether \(P_\beta\) is true:
this last contradiction uses a rejection of
a word in \(L\), not necessarily a new insertion
at that particular stage.

The cases exhaust all possibilities. In each
possible case the construction is a primitive-recursive
text of a target in \(\mathcal C\) on which
the normalized learner fails. By (PR.2), its
original total computable learner fails too.
This proves the asserted tester obstruction,
without a semantic hypothesis-comparison oracle
or a totality test for intermediate program guesses.
\end{proof}
\end{annexitem}
